\chapter{Elements, style, colour, border}
\label{ch:elements}

\code{.build()} converts the compile-time tree from
Chapter~\ref{ch:dsl} into an \code{Element}. That \code{Element} is
the input to the layout engine, the renderer, and every widget. This
chapter dissects \code{Element} and the value types it depends on.

\begin{aside}
Chapter~\ref{ch:dsl} was about the compile-time side --- the
elaborate type machinery that makes \code{t<"hi"> | Bold} legal and
\code{v(\dots) | bcol<\dots>} (without a border) illegal. This
chapter is about the runtime side: the small, boring, fast data
structure the renderer actually consumes. The contrast is part of
the trick. The compile-time DSL is intricate because we want the
compiler to do interesting work; the runtime representation is
plain because we want the renderer to be fast.
\end{aside}

\section{The runtime tree, in one sentence}
\label{sec:elem:sentence}

Everything in this chapter rotates around one idea: an
\textbf{element tree} is a tree where every node is one of three
things.

\begin{itemize}[leftmargin=*]
  \item \textbf{Text} --- a leaf, holding bytes and styling.
  \item \textbf{Box} --- a container, holding a layout configuration
    and a list of child elements.
  \item \textbf{List} --- a transparent grouping, ``these elements
    belong to my parent's children'', used for splicing.
\end{itemize}

Nothing more. Every shape you have seen \maya{} produce ---
dashboards, list views, scrolling chats, modal dialogs --- is built
out of those three nodes nested arbitrarily.

\section{What is a tagged union?}
\label{sec:elem:tagged-union}

Before we go into the variant: a small piece of language theory you
will use over and over in this book.

A \textbf{tagged union} (also called a \emph{sum type} or
\emph{discriminated union}) is a value that is \emph{one of} a fixed
set of alternatives, with a tag indicating which one. The phrase
``one of'' matters: at any moment, the value is exactly one
alternative, never two, never zero. The compiler can check that you
have handled every case.

Contrast with the other big composite type:

\begin{itemize}[leftmargin=*]
  \item A \textbf{product type} (a struct, class, tuple) is
    \emph{this and this and this}. A \code{Point\{int x, int y\}}
    has both an \code{x} \emph{and} a \code{y}, always.
  \item A \textbf{sum type} (a tagged union, variant) is
    \emph{this or this or this}. A \code{Result} is either an
    \code{Ok} or an \code{Error}, never both at once.
\end{itemize}

Natural-language analogies break almost immediately when you push
them, but here is a small one: a parking space's state is a sum
type (\emph{Empty} or \emph{Occupied}, never both); a car's state
is a product type (\emph{has colour} \textbf{and} \emph{has
licence plate}, both at once).

In modern \cpp{} the sum-type primitive is \code{std::variant<A, B,
C>}, which holds exactly one of \code{A}, \code{B}, or \code{C} and
remembers which. You inspect it with \code{std::visit}, which calls
your visitor with the actual stored alternative. The compiler
verifies that your visitor handles every case --- forget one and
you get an error.

\begin{lstlisting}
#include <variant>
using Pet = std::variant<Dog, Cat, Hamster>;

std::string describe(const Pet& p) {
    return std::visit(overload{
        [](const Dog& d)     { return "a good boy named " + d.name; },
        [](const Cat& c)     { return c.name + " who is judging you"; },
        [](const Hamster& h) { return "a hamster named " + h.name; },
    }, p);
}
\end{lstlisting}

Three branches, three alternatives. Add a fourth alternative
(\code{Goldfish}) and the visitor stops compiling --- the compiler
lists the missing branch by name.

\begin{idea}
A sum type is the data-structures version of an \code{if}/\code{else}
chain. It is \emph{closed} (the alternatives are listed in the
type), \emph{exhaustive} (the compiler nags you about missing
cases), and \emph{cheap} (one word for the tag, plus enough storage
for the largest alternative).
\end{idea}

\subsection{Why \maya{} uses sum types everywhere}

Look through the source and you will find sum types in almost every
file:

\begin{itemize}[leftmargin=*]
  \item \code{Element = variant<BoxElement, TextElement, ElementList>}
    --- this chapter.
  \item \code{Cmd<Msg>::Variant} --- algebra of effects
    (Chapter~\ref{ch:elm}).
  \item \code{Sub<Msg>::Variant} --- algebra of subscriptions
    (Chapter~\ref{ch:elm}).
  \item \code{Event = variant<KeyEvent, MouseEvent, ResizeEvent, PasteEvent>}
    --- input events (Chapter~\ref{ch:events}).
  \item \code{RenderOp = variant<EmitText, MoveCursor, SetStyle, ...>}
    --- renderer output (Chapter~\ref{ch:ansi}).
  \item \code{expected<T, E>} = \code{variant<T, E>} essentially ---
    success or failure (Chapter~\ref{ch:cpp}).
\end{itemize}

The reason is uniform: each of these is a finite, named set of
possibilities. Encoding them as sum types gives you exhaustiveness
checking and removes whole classes of \code{nullptr}-ish bugs. When
you learn to recognise the pattern, you start seeing it everywhere
in well-designed C++ codebases.

\section{\code{Element}: a tagged union}

\code{Element} is a \code{std::variant} of three alternatives:

\begin{lstlisting}
using Element = std::variant<BoxElement, TextElement, ElementList>;
\end{lstlisting}

\begin{itemize}[leftmargin=*]
  \item \code{TextElement} --- a piece of styled text, possibly with
    a wrap policy.
  \item \code{BoxElement} --- a container with flex layout, padding,
    border, gap, etc., and a vector of child elements.
  \item \code{ElementList} --- a flat sequence of elements that
    should be inlined into the parent's children (used by \code{map}).
\end{itemize}

Three alternatives, a closed set. Every algorithm in \maya{} that
walks an element tree (layout, render, hit-testing, search) ends up
in a \code{std::visit} on this variant. Add a fourth alternative
tomorrow and the compiler tells you all the places to update.

\subsection{\code{TextElement}}

\begin{lstlisting}
struct TextElement {
    std::string text;        // the bytes (UTF-8)
    Style       style;       // colours + attributes
    TextWrap    wrap = TextWrap::Wrap;
    HyperlinkId link{};      // optional OSC-8 link
};
\end{lstlisting}

A leaf in the tree. The string is owned (so the source data can be a
temporary). Styling lives in a \code{Style} struct (next section).
\code{TextWrap} is one of \code{Wrap}, \code{NoWrap}, \code{Truncate},
\code{Ellipsize}.

\subsection{\code{BoxElement}}

\begin{lstlisting}
struct BoxElement {
    layout::FlexStyle    style;     // direction, padding, gap, ...
    std::optional<Border> border;
    std::vector<Element>  children;
    std::optional<std::string> title;  // border title
    // ... a few more renderer hints
};
\end{lstlisting}

The composite. \code{style} carries every layout property as integer
cell counts; \code{border} is set when the DSL chain included
\code{border\_<...>}. \code{title} is set when a border title was
piped on.

\subsection{\code{ElementList}}

\begin{lstlisting}
struct ElementList { std::vector<Element> items; };
\end{lstlisting}

A pseudo-element. The renderer treats it as ``these children belong
to my parent.'' It is the runtime equivalent of a parameter pack
expansion.

\section{\code{Style}}

\code{Style} is a small POD struct describing the appearance of a run
of text:

\begin{lstlisting}
struct Style {
    Color     fg{};
    Color     bg{};
    Attribute attrs{};   // bold | italic | underline | dim | ...

    constexpr Style with_bold()      const noexcept;
    constexpr Style with_italic()    const noexcept;
    constexpr Style with_fg(Color c) const noexcept;
    constexpr Style with_bg(Color c) const noexcept;
    constexpr Style merge(Style other) const noexcept;
    // ...
};
\end{lstlisting}

Two flavours exist:

\begin{itemize}[leftmargin=*]
  \item \code{Style} (runtime), used in \code{TextElement}.
  \item \code{CTStyle} (compile-time), used as an NTTP in
    \code{TextNode}. It is structurally identical but lives in
    template parameters and is constructed with \code{constexpr}
    factory functions (\code{CTStyle::bold()},
    \code{CTStyle::fg(R, G, B)}, etc.).
\end{itemize}

The DSL pipes operate on \code{CTStyle}; \code{.build()} calls
\code{CTStyle.runtime()} to project back to a \code{Style}. The
distinction matters because NTTPs require structural types whose
fields are compared by value, while runtime \code{Style} can have
fancier constructors.

\subsection{\code{Attribute} bitset}

Attributes (bold, dim, italic, underline, strikethrough, inverse,
blink) are a bitset:

\begin{lstlisting}
enum class Attribute : std::uint8_t {
    None       = 0,
    Bold       = 1 << 0,
    Dim        = 1 << 1,
    Italic     = 1 << 2,
    Underline  = 1 << 3,
    Strike     = 1 << 4,
    Inverse    = 1 << 5,
    Blink      = 1 << 6,
};

constexpr Attribute operator|(Attribute a, Attribute b) noexcept;
constexpr bool      has(Attribute set, Attribute bit) noexcept;
\end{lstlisting}

Combining attributes is bitwise OR. Checking is bitwise AND. The
renderer translates each bit to its SGR parameter (1 for bold, 2 for
dim, 3 for italic, etc.) when serialising.

\section{\code{Color}}

A \code{Color} is a 24-bit RGB triple plus a tag for ``default''
(unspecified, inherit from terminal):

\begin{lstlisting}
struct Color {
    std::uint8_t r = 0, g = 0, b = 0;
    bool         is_default = true;   // means "use terminal default"
};

constexpr Color rgb(uint8_t r, uint8_t g, uint8_t b) noexcept {
    return {r, g, b, false};
}
\end{lstlisting}

\maya{} only deals in 24-bit truecolour internally. The 256-colour and
16-colour palettes were popular when terminals had limited support;
modern terminals all do truecolour. If you really need 256-colour
output (because you are targeting a niche emulator), the writer can
quantise on the way out.

\subsection{Themes}

A \textbf{theme} is a 24-slot palette of colours --- one for primary,
secondary, success, warning, error, plus background tiers and text
tiers. \maya{}'s built-in themes (\code{dark}, \code{light},
\code{terminal}) are \code{constexpr} aggregates; you can derive your
own and pass it in \code{RunConfig}.

\begin{lstlisting}
struct Theme {
    Color bg_primary;
    Color bg_secondary;
    Color text_primary;
    Color text_secondary;
    Color accent;
    Color success;
    Color warning;
    Color error;
    // ... 24 slots total
};
\end{lstlisting}

The \code{Ctx} parameter to a render function (in the quick-tool
\code{run(...)}) carries the active theme. Widgets read it instead of
hard-coding colours --- so the same widget tree looks right under
dark and light themes.

\section{\code{Border}}

A border is described by its style (single, double, round, thick,
classic, arrow), its colour (or default), and which sides are drawn:

\begin{lstlisting}
enum class BorderStyle : std::uint8_t {
    None, Single, Double, Round, Bold, Classic, Arrow,
};

struct BorderSides {
    bool top = true, right = true, bottom = true, left = true;

    static constexpr BorderSides all()        { return {true, true, true, true}; }
    static constexpr BorderSides none()       { return {false,false,false,false}; }
    static constexpr BorderSides horizontal() { return {true,false,true,false}; }
    static constexpr BorderSides vertical()   { return {false,true,false,true}; }
};

struct Border {
    BorderStyle style = BorderStyle::Single;
    Color       colour{};
    BorderSides sides = BorderSides::all();
};
\end{lstlisting}

Each \code{BorderStyle} maps to a set of eight glyphs (corners + edges)
that the renderer paints. The mapping is a \code{constexpr}
\code{std::array} in \file{maya/style/border.hpp}.

\subsection{\code{Sep} and \code{VSep}}

The DSL's \code{sep} and \code{vsep} nodes are syntactic sugar for a
box with a one-side-only border --- horizontal or vertical. The same
machinery; the renderer treats them identically.

\section{\code{BoxBuilder}: runtime construction}

The DSL \code{.build()} fans out into a \code{BoxBuilder} which
assembles a \code{BoxElement}:

\begin{lstlisting}
class BoxBuilder {
public:
    BoxBuilder(FlexDirection dir, BoxCfg cfg)
      : dir_(dir), cfg_(cfg) {}

    void add(Element child) { children_.push_back(std::move(child)); }
    void add(ElementList list) {
        children_.insert(children_.end(),
                         std::make_move_iterator(list.items.begin()),
                         std::make_move_iterator(list.items.end()));
    }

    Element finalize() &&;
private:
    FlexDirection         dir_;
    BoxCfg                cfg_;
    std::vector<Element>  children_;
};
\end{lstlisting}

Two overloads of \code{add}: one for ordinary children, one for an
\code{ElementList} which is spliced. \code{finalize()} translates
\code{BoxCfg} into a \code{layout::FlexStyle}, attaches a \code{Border}
if requested, wraps it in a \code{BoxElement}, and returns the
variant.

\section{Walking an element tree}

Every algorithm that processes an element tree is a recursive
\code{std::visit}. Here is the shape of a pretty-printer:

\begin{lstlisting}
void print(const Element& e, int depth = 0) {
    std::string indent(depth * 2, ' ');
    std::visit(overload{
        [&](const TextElement& t) {
            std::println("{}TEXT  : \"{}\"", indent, t.text);
        },
        [&](const BoxElement& b) {
            std::println("{}BOX   : dir={}, {} children",
                         indent, int(b.style.direction), b.children.size());
            for (const auto& c : b.children) print(c, depth + 1);
        },
        [&](const ElementList& l) {
            for (const auto& c : l.items) print(c, depth);
        }
    }, e);
}
\end{lstlisting}

This is the template every renderer-side pass follows: visit, branch
on alternative, recurse on children.

\section{Hyperlinks}

A small but neat feature: text elements can carry a hyperlink. The
renderer emits an OSC-8 sequence around the run:

\begin{lstlisting}
ESC ] 8 ; ; URL ESC \  text  ESC ] 8 ; ; ESC \
\end{lstlisting}

Modern emulators turn that into a clickable link. \maya{} interns the
URL into a \code{HyperlinkId} so the per-cell representation stays
compact (more on that in Chapter~\ref{ch:canvas}).

\section{Why so many small structs?}

A reasonable question: why isn't \code{Element} a class with virtual
methods? \maya{}'s answer: virtual dispatch precludes a packed cell
representation, prevents \code{constexpr} construction, and makes
serialization harder. The variant + visitor pattern gives you
polymorphism without the overhead --- \code{std::visit} compiles to a
jump table or a chain of branches, both faster than a virtual call,
and the data layout stays cache-friendly.

\section{The forward-declaration dance}
\label{sec:elem:forward-decl}

If you open \file{maya/include/maya/element/element.hpp} you will see
an unusual little ritual: \code{Element} is forward-declared,
\code{box.hpp} and \code{text.hpp} are then included, and only after
that is \code{Element} actually defined. The comments in the file
explain why; the explanation is worth understanding because the same
problem will come up in your own code.

The shape is:

\begin{lstlisting}
// Step 1: forward-declare Element
struct Element;
struct ElementList;

// Step 2: include the concrete element headers
#include "box.hpp"      // BoxElement, which has std::vector<Element>
#include "text.hpp"     // TextElement

// Step 3-4: now define Element and ElementList
struct ElementList { std::vector<Element> items; ... };
struct Element     { /* the variant */ };
\end{lstlisting}

Why this order? Because \code{BoxElement} contains
\code{std::vector<Element>}, which means \code{BoxElement}'s
definition needs to know about \code{Element}. But \code{Element} is
the variant of \code{BoxElement}, \code{TextElement}, and
\code{ElementList}, so its definition needs to know about
\code{BoxElement} --- a circular dependency.

\cpp{} resolves this with \textbf{forward declarations}: a
forward-declared type is enough for the compiler to know
``\code{Element} is a type, you can have pointers and references to
it, and \code{std::vector<Element>} works because the vector holds
an indirect pointer to its element type internally''. The full
definition is supplied later, after the dependent pieces are in
place.

The technique generalises: any time you have a tree where a node's
field holds a vector / optional / pointer of the same type, you
use forward declaration to break the cycle.

\begin{idea}
A tree type that contains itself in a field is a circular type
definition. \cpp{} breaks the circle with a forward declaration:
the outer type is named first, the indirect field can then refer to
it, and the full definition follows.
\end{idea}

\section{Constructing an \code{Element} by hand}
\label{sec:elem:by-hand}

The DSL hides almost all element construction. For tests,
diagnostics, and writing widget primitives that work below the DSL
layer, you sometimes need to make an \code{Element} directly.

The direct way:

\begin{lstlisting}
#include <maya/element/element.hpp>

using namespace maya;

// A leaf text element
Element title = TextElement{
    .text  = "Hello",
    .style = Style{}.with_bold().with_fg(Color::rgb(100, 180, 255)),
};

// A box containing two children
Element panel = BoxElement{
    .style    = layout::FlexStyle{ .direction = FlexDirection::Column,
                                   .padding   = {.top=1, .right=2,
                                                 .bottom=1, .left=2} },
    .border   = Border{.style = BorderStyle::Round},
    .children = {
        std::move(title),
        Element{TextElement{.text = "world"}}
    },
};
\end{lstlisting}

The braces in the field initialisers are designated initialisers
(\cpp{}20); they make the field names explicit so the code reads
like data. The whole construction is value-based: no \code{new},
no \code{shared\_ptr}, just nested structs.

The more idiomatic way uses the \code{box()} builder:

\begin{lstlisting}
#include <maya/element/builder.hpp>

Element panel = maya::detail::box()
    .direction(FlexDirection::Column)
    .padding(1, 2, 1, 2)
    .border(BorderStyle::Round)
    .border_color(Color::rgb(60, 65, 80))
    (text_el("Hello").bold().fg(100, 180, 255),
     text_el("world"));
\end{lstlisting}

The builder is a fluent chain that mirrors the DSL pipes; its final
\code{operator()} takes the children. This is what the DSL's
\code{.build()} ultimately calls --- the builder lives in
\file{maya/include/maya/element/builder.hpp}.

\section{The colour model up close}
\label{sec:elem:color}

We established that \code{Color} is a 24-bit RGB triple plus an
``is default'' bit. Worth pausing on what that bit costs and what it
gains.

Without it, you would need a sentinel: \code{Color\{0, 0, 0\}}
meaning either ``actual black'' or ``unset''. The runtime would
have no way to tell ``the user explicitly chose black'' from ``the
user hasn't specified'' --- and the renderer needs to know, because
the terminal's default foreground colour may not be black (it is
white-on-black for dark themes, dark-on-white for light themes).

With \code{is\_default = true}, the renderer can emit
\code{ESC[39m} (\emph{reset foreground to terminal default}) rather
than a specific colour. The cell inherits whatever the user's
terminal preferences say.

\begin{lstlisting}
Color c1 = Color::rgb(0, 0, 0);                // explicit black
Color c2 = Color::default_();                  // terminal default
Color c3{};                                    // also terminal default
// c1.is_default == false, c2.is_default == true
\end{lstlisting}

The difference matters when a \maya{} program runs in someone
else's terminal: their carefully-chosen colour scheme is preserved
for anything you did not explicitly style.

\subsection{The hex shorthand}

Designers think in hex; the runtime has a constructor for that:

\begin{lstlisting}
Color accent = Color::hex(0x64B4FF);    // = rgb(100, 180, 255)
Color danger = Color::hex(0xFF4444);    // = rgb(255, 68, 68)
\end{lstlisting}

The DSL has matching \code{fg<0x\dots>} and \code{bg<0x\dots>}
pipes (Chapter~\ref{ch:dsl}) which call the same constructor at
compile time.

\subsection{Truecolour, 256-colour, 16-colour}

A brief detour. The terminal world has three historical colour
palettes:

\begin{itemize}[leftmargin=*]
  \item \textbf{8 / 16 colours.} The classic ANSI palette. SGR codes
    30--37 (foreground), 40--47 (background), 90--97 (bright
    foreground), 100--107 (bright background). Universally
    supported.
  \item \textbf{256 colours.} An xterm extension; a 6×6×6 colour
    cube plus 24 greys plus the original 16. Encoded as
    \code{ESC[38;5;<n>m}. Most emulators support it.
  \item \textbf{Truecolour (24-bit).} Three bytes per channel.
    Encoded as \code{ESC[38;2;<R>;<G>;<B>m}. Modern emulators ---
    every one made in the last decade --- support this.
\end{itemize}

\maya{} stores colours internally as truecolour and lets the
renderer worry about quantisation if the target terminal does not
support it. In practice, ``the target terminal does not support
it'' is so rare you can ignore it; if you really need to target a
2003 emulator over a hardened serial link, set
\code{RunConfig::color\_mode = ColorMode::ANSI16} and a quantiser
will snap each truecolour to the nearest ANSI colour on the wire.

\begin{funfact}
24-bit truecolour over a terminal was proposed in the early 2000s
and snuck into emulators piecemeal over the next decade. The
official standard governing it (\code{ITU-T T.416}) was published
in 1993 but ignored for fifteen years. Today, your
\code{printf '\textbackslash{}e[38;2;255;128;0mhi'} works in every
mainstream emulator without anyone needing to think about it. The
gap between a feature existing and it being usable is sometimes
decades long.
\end{funfact}

\section{The border gallery}
\label{sec:elem:border-gallery}

Each \code{BorderStyle} maps to a set of glyphs the renderer paints
at the cells immediately inside the box's bounding rectangle. The
styles are not interchangeable; the glyph choice signals affect.

\paragraph{\code{None}.} No border, no glyphs. The box's content
fills its rectangle.

\paragraph{\code{Single}.} The default. Light-weight box-drawing
characters from Unicode's \code{U+2500}--\code{U+257F} block:
horizontal bar (\code{U+2500}), vertical bar (\code{U+2502}),
top-left corner (\code{U+250C}), and so on.

\paragraph{\code{Round}.} Same edges as \code{Single}, but the four
corners use the rounded variants (\code{U+256D}, \code{U+256E},
\code{U+256F}, \code{U+2570}). Looks friendlier on most fonts.

\paragraph{\code{Double}.} Double-line glyphs (\code{U+2550} for
horizontal, \code{U+2551} for vertical, \code{U+2554} for top-left).
The classic \emph{IBM} look from DOS-era applications.

\paragraph{\code{Bold} (alias \code{Thick}).} Heavy box-drawing
glyphs (\code{U+2501}, \code{U+2503}, \code{U+250F}). Useful for
emphasis or selected state.

\paragraph{\code{Classic}.} ASCII-only fallback using \code{+},
\code{-}, \code{|}. Use for legacy terminals that do not render
Unicode box-drawing characters.

\paragraph{\code{Arrow}.} A novelty style with arrow heads pointing
into the corners. Good for ``selected'' or ``focus'' indicators on a
subpanel.

The glyph table for every style lives in
\file{maya/include/maya/style/border.hpp} as a \code{constexpr}
\code{std::array}, one row per style. You can read all of it in one
minute.

\subsection{Partial-side borders}

\code{BorderSides} controls which of the four edges actually paint.
The defaults are all four; the helpers \code{horizontal()} and
\code{vertical()} pick top+bottom or left+right respectively. This
is how \code{sep} and \code{vsep} are implemented --- they are
boxes with one side worth of border, drawn as a thin separator.

\begin{lstlisting}
// A separator under a heading
auto heading = v(
    t<"Section"> | Bold,
    sep,                       // = a row with .border.sides = horizontal()
    body
);
\end{lstlisting}

A separator costs one row of height. Vertical separators (\code{vsep})
cost one column of width and appear inside row layouts.

\section{Patterns for walking the tree}
\label{sec:elem:walk-patterns}

The pretty-printer earlier was the simplest possible tree walk. Real
algorithms in \maya{} have a few additional patterns; learning them
lets you build your own analyses (linters, accessibility checks,
layout debuggers).

\subsection{Counting nodes}

\begin{lstlisting}
int count(const Element& e) {
    return std::visit(overload{
        [](const TextElement&)    { return 1; },
        [&](const BoxElement& b)  {
            int n = 1;
            for (const auto& c : b.children) n += count(c);
            return n;
        },
        [&](const ElementList& l) {
            int n = 0;
            for (const auto& c : l.items) n += count(c);
            return n;
        },
    }, e);
}
\end{lstlisting}

\code{ElementList} does not count itself (it is a transparent
fragment), but its items do. \code{BoxElement} counts itself plus
its children.

\subsection{Finding all text}

Flatten an \code{Element} into a vector of strings (useful for
testing, search, screen readers):

\begin{lstlisting}
void texts(const Element& e, std::vector<std::string>& out) {
    std::visit(overload{
        [&](const TextElement& t) { out.push_back(t.text); },
        [&](const BoxElement& b)  {
            for (const auto& c : b.children) texts(c, out);
        },
        [&](const ElementList& l) {
            for (const auto& c : l.items) texts(c, out);
        },
    }, e);
}
\end{lstlisting}

The same scaffold; the only thing that changes is what you do at
each node.

\subsection{Transforming the tree}

What if you want to return a \emph{new} tree, with every text node
uppercased?

\begin{lstlisting}
Element upper(const Element& e) {
    return std::visit(overload{
        [](const TextElement& t) -> Element {
            auto copy = t;
            std::transform(copy.text.begin(), copy.text.end(),
                           copy.text.begin(), ::toupper);
            return copy;
        },
        [&](const BoxElement& b) -> Element {
            BoxElement nb = b;
            for (auto& c : nb.children) c = upper(c);
            return nb;
        },
        [&](const ElementList& l) -> Element {
            ElementList nl = l;
            for (auto& c : nl.items) c = upper(c);
            return nl;
        },
    }, e);
}
\end{lstlisting}

The shape is the same again --- visit, recurse on children, build
the new node. With trees of plain values, transformation is
non-destructive and trivially correct: the original tree is
untouched.

\subsection{Hit testing}

Given a screen coordinate, which \code{Element} is at that position?
The layout engine (Chapter~\ref{ch:layout}) assigns each node a
rectangle; hit testing walks the tree picking the deepest node
whose rectangle contains the point. The pattern is identical:
\code{std::visit}, recurse, return.

\begin{idea}
Every tree pass in \maya{} is \code{std::visit} + recursion. The
variant is closed, the visitor is exhaustive, and adding a new node
kind would force every visitor in the codebase to add a branch ---
in one compiler error message, with one source location each. This
is why the variant pattern scales.
\end{idea}

\section{The \code{Style} merge rule}
\label{sec:elem:style-merge}

A detail that catches people out: how does \code{Style} composition
actually work? Two styles, one inheriting from a parent and one
applied locally, must combine somehow.

The rule, implemented by \code{Style::merge}:

\begin{itemize}[leftmargin=*]
  \item Colours: the second style's \code{fg} / \code{bg} replace
    the first's \emph{only if they are not default}. The default
    sentinel acts as ``do not override''.
  \item Attributes: a bitwise OR. Bold can be added on top of dim;
    underline can be added on top of bold. There is no ``unbold''
    attribute in the merged form --- styles only accumulate.
\end{itemize}

In practice this means:

\begin{lstlisting}
Style base = Style{}.with_fg(Color::rgb(100, 100, 100));
Style bold = Style{}.with_bold();
Style merged = base.merge(bold);
// merged.fg    == rgb(100, 100, 100)
// merged.attrs == Bold
\end{lstlisting}

The parent's grey foreground is preserved (bold did not specify a
colour), and the bold attribute is added. If you later wanted to
actively reset bold, you would do it by constructing a fresh style,
not by ``negating'' in the merge.

\section{Hyperlinks}

A small but neat feature: text elements can carry a hyperlink. The
renderer emits an OSC-8 sequence around the run:

\begin{lstlisting}
ESC ] 8 ; ; URL ESC \  text  ESC ] 8 ; ; ESC \
\end{lstlisting}

Modern emulators turn that into a clickable link. \maya{} interns the
URL into a \code{HyperlinkId} so the per-cell representation stays
compact (more on that in Chapter~\ref{ch:canvas}).

\begin{lstlisting}
Element docs = TextElement{
    .text = "see docs",
    .style = Style{}.with_underline(),
    .link  = intern_hyperlink("https://example.org/docs"),
};
\end{lstlisting}

The \code{intern\_hyperlink} call returns a small integer; the URL
string lives in a per-frame interning table, freed at end of frame.
This keeps the per-cell representation a small fixed-size struct
rather than a string.

\section{Pitfalls and anti-patterns}
\label{sec:elem:pitfalls}

\begin{warning}
\textbf{Storing an \code{Element} across frames.} The element tree is
built fresh on every frame; storing a node from frame \code{N} for
use in frame \code{N+1} is not catastrophic, but it defeats the
rebuild-from-model discipline and is usually a sign your model is
missing a field. Put the data in the model, build the element from
it.
\end{warning}

\begin{warning}
\textbf{Mutating an \code{Element} after construction.} The runtime
may assume the element you returned from \code{view()} is
immutable for the duration of the frame. Modifying it after the
fact (from a background thread, say) is undefined behaviour. If you
need to ``change the UI'' from a background thread, send a
\code{Msg} through \code{Cmd::task}'s dispatch; \code{view} will be
re-invoked with the new model.
\end{warning}

\begin{warning}
\textbf{Using raw heap allocation.} \code{new BoxElement\{\dots\}}
and returning a pointer breaks the value-semantics invariant the
rest of the library assumes. The \code{Element} variant is meant to
be constructed by value and moved; mixing pointer ownership in
will break copy / move semantics in confusing ways.
\end{warning}

\begin{warning}
\textbf{Building deeply nested boxes when a flat list would do.}
Layout time grows roughly linearly with node count. A list of 100
rows nested inside a vertical box is fine; a list of 100 rows where
each row is a box containing a box containing a box of cells is
slow. Keep the tree as flat as you can.
\end{warning}

\begin{recap}
The runtime tree is \code{Element = variant<BoxElement, TextElement,
ElementList>}. Style, Color, and Border are small POD structs;
attributes are a bitset; themes are 24-slot palettes. Every
algorithm in \maya{} walks the tree with \code{std::visit}.
\code{BoxBuilder} bridges the compile-time DSL to the runtime tree.
\end{recap}

\section*{Workshop}

\begin{enumerate}[leftmargin=*]
  \item Open \file{maya/include/maya/element/element.hpp} and find
    the \code{Element} alias. Note that it is exactly the variant
    described in this chapter.
  \item Write a function
    \code{int count\_text\_nodes(const Element\&)} using
    \code{std::visit}. Test it on a panel from
    Chapter~\ref{ch:dsl}.
  \item Construct an \code{Element} by hand --- not via the DSL ---
    using \code{BoxBuilder}. Render it with \code{print(...)}.
\end{enumerate}

\section{Sum types and the shape of UI data}

Before we tour every variant in \maya{}'s element tree, it is
worth pausing on \emph{why} variants are the right tool for UI
trees. The decision was not aesthetic; it was forced by what UI
trees look like.

\subsection*{Product vs sum types}

The two basic shapes of compound data are products and sums:

\begin{itemize}[leftmargin=*]
  \item A \textbf{product type} carries one of each member. A
    \code{struct Point \{ int x; int y; \};} is a product: it has
    both an x \emph{and} a y.
  \item A \textbf{sum type} carries one of several possibilities.
    A \code{std::variant<int, double, std::string>} is a sum: it
    is an int \emph{or} a double \emph{or} a string, never more
    than one at a time.
\end{itemize}

A UI tree node is fundamentally a sum: a node is \emph{either} a
box \emph{or} a piece of text \emph{or} a list of children. It is
never simultaneously a box and a text. The right C++ tool for ``X
or Y or Z'' is \code{std::variant}.

\subsection*{Why not inheritance?}

For decades the conventional answer was a class hierarchy:

\begin{lstlisting}
class Element { /* virtual ~Element() = 0; */ };
class Box  : public Element { ... };
class Text : public Element { ... };
// hand around Element* everywhere
\end{lstlisting}

This works, but it pays for flexibility you do not need:

\begin{itemize}[leftmargin=*]
  \item \textbf{Heap allocation.} Polymorphism through base
    pointers forces heap storage; an \code{Element*} cannot be a
    stack value of unknown size.
  \item \textbf{Indirection.} Every method call is a virtual
    dispatch through a vtable. Cache-cold for typical UI trees.
  \item \textbf{Open set.} Anyone can derive a new
    \code{Element}. That sounds like a feature; for a library
    that wants to optimise diff and layout, it is a curse. The
    set of node kinds must be closed for the compiler to
    exhaustively check switches.
  \item \textbf{Object lifetime.} \code{std::unique\_ptr<Element>}
    works, but every operation that wants to copy a tree has to
    deep-clone via a virtual \code{clone()}.
\end{itemize}

A closed sum type fixes all four. The set of kinds is known at
compile time; storage is inline; dispatch is a small switch the
optimiser inlines; copies are bytewise.

\subsection*{std::variant: how it actually works}

A \code{std::variant<A, B, C>} is, internally, a tagged union: a
small integer (the \emph{index}) plus enough storage for the
largest alternative, aligned for any of them.

\begin{lstlisting}
// Simplified internal layout.
struct VariantABC {
    std::size_t index;        // 0, 1, or 2
    union {
        A a;
        B b;
        C c;
    };
};
\end{lstlisting}

When you write \code{std::get<A>(v)}, it checks \code{index == 0}
and returns a reference to the active member. When you write
\code{std::visit(f, v)}, the implementation generates a switch
on \code{index} that calls \code{f} with the active member.

In release builds the optimiser usually inlines the switch into
the caller. The result is dispatch as fast as a virtual call
\emph{and} no indirection through the heap.

\begin{funfact}
\code{std::variant} was added in C++17. Before that, the
C++ ecosystem had Boost.Variant (since 2002) and a handful of
hand-rolled tagged unions in every large codebase. The standard
finally caught up to ML, which had had tagged unions since 1973.
\end{funfact}

\subsection*{std::visit and the overload trick}

\code{std::visit(f, v)} requires \code{f} to be callable on every
alternative. The cleanest way to write \code{f} is the
\emph{overload pattern}:

\begin{lstlisting}
template <class... Fs> struct overload : Fs... { using Fs::operator()...; };
template <class... Fs> overload(Fs...) -> overload<Fs...>;

// Usage:
auto result = std::visit(overload{
    [](Box b)  { return "box"; },
    [](Text t) { return "text"; },
    [](auto)   { return "?"; }
}, element);
\end{lstlisting}

The \code{overload} struct inherits from a pack of lambdas and
uses \code{using Fs::operator()...} to bring them all into scope
as overloads of a single \code{operator()}. The compiler picks
the one whose parameter matches the active alternative.

This pattern is so common that C++26 is considering builtin
support. For now, every \maya{} translation unit has one
\code{overload} struct.

\begin{recap}
A tagged union is a value-typed sum: index + storage. Visiting it
is a switch the compiler optimises into a jump table. Inheritance
is the wrong tool for UI nodes because the set of kinds is
closed; sums make the closure explicit and let the compiler
check exhaustiveness.
\end{recap}

\section{Every Element variant, by kind}

With the theory settled, here is the actual variant in
\file{maya/include/maya/element/element.hpp}:

\begin{lstlisting}
using Element = std::variant<BoxElement, TextElement, ElementList>;
\end{lstlisting}

Three alternatives. Two represent a node, one represents a
container (a flat list of nodes, used where layout does not
benefit from a box). Each is a struct with public fields.

\subsection*{TextElement}

A single run of text with one style:

\begin{lstlisting}
struct TextElement {
    std::variant<std::string_view, std::string> text;
    Style style;
    std::optional<std::string> hyperlink;
};
\end{lstlisting}

The \code{text} field is itself a small variant: a non-owning view
(for compile-time literals) or an owning string (for runtime
content). This is the runtime echo of the DSL's TextNode /
RuntimeTextNode split.

\code{style} carries colours, attributes, and the foreground/
background tags. \code{hyperlink} is the OSC 8 target URL, if any.

\subsection*{BoxElement}

A container with optional decoration:

\begin{lstlisting}
struct BoxElement {
    std::vector<Element> children;
    FlexStyle layout;
    std::optional<Border> border;
    Padding padding;
    Style style;
    Visibility visibility;
};
\end{lstlisting}

This is the workhorse. \code{children} is a flat list of any
element kind. \code{layout} contains the flexbox properties
(direction, justify, align, grow, shrink) covered in
Chapter~\ref{ch:layout}. \code{border} is an optional decoration
struct; \code{padding} is per-side integer cells; \code{style}
applies to the entire box content.

\code{visibility} is a small enum: \code{Shown}, \code{Hidden}
(layout reserves space, nothing drawn), \code{Collapsed} (no
space, no drawing). The distinction matters when you want a
toggle that does not jiggle the layout.

\subsection*{ElementList}

A list of elements without a box around them:

\begin{lstlisting}
struct ElementList {
    std::vector<Element> items;
};
\end{lstlisting}

Useful when you want to return ``several elements'' from a helper
function without forcing them into a box. The layout pass treats
an \code{ElementList} as if its items had been spliced in place.

\section{Forward declarations: the chicken-and-egg dance}

A \code{BoxElement} contains \code{std::vector<Element>}. An
\code{Element} is a \code{std::variant<BoxElement, \ldots>}.
This is mutually recursive: \code{Element} depends on
\code{BoxElement}, and \code{BoxElement} depends on
\code{Element}.

C++ does not allow circular type definitions. You cannot write:

\begin{lstlisting}
struct BoxElement { std::vector<Element> children; };  // error: Element undefined
using Element = std::variant<BoxElement, TextElement>;
\end{lstlisting}

The standard library solves it in two steps:

\begin{enumerate}[leftmargin=*]
  \item \code{std::vector<T>} requires only that \code{T} be a
    forward-declared type, not a complete one. (This is unusual
    for the standard library --- most containers require complete
    types --- but vector got the relaxed rule in C++17.)
  \item Therefore you can declare \code{BoxElement} containing
    \code{std::vector<Element>} as long as \code{Element} is at
    least forward-declared at that point.
\end{enumerate}

The pattern in \maya{}:

\begin{lstlisting}
// 1. Forward-declare Element.
struct Element;

// 2. Define the alternatives. They may contain vectors of Element.
struct TextElement { /* ... */ };
struct BoxElement  { std::vector<Element> children; /* ... */ };
struct ElementList { std::vector<Element> items; };

// 3. Now define Element itself.
struct Element {
    std::variant<BoxElement, TextElement, ElementList> v;
};
\end{lstlisting}

In practice \maya{} wraps the variant in a struct (rather than
using a \code{using Element = std::variant<\ldots>;}) so it can
add convenience methods like \code{is\_text()} and
\code{as\_box()}.

\begin{aside}
This dance generalises. Any recursive data structure in C++ uses
the same trick: forward-declare the outer type, define the inner
types that contain a vector or pointer to the outer type, then
define the outer type. The classic example is a JSON value, where
\code{Object} maps strings to values and \code{Array} is a vector
of values.
\end{aside}

\section{Constructing an Element by hand}

The DSL builds elements for you, but sometimes you need to
assemble one without the compile-time machinery (loaded layouts,
user-themable trees, debugging). C++26's designated initialisers
make this readable:

\begin{lstlisting}
Element greeting{
    .v = TextElement{
        .text  = std::string_view{"Hello, maya"},
        .style = Style{
            .fg = Color::rgb(220, 220, 240),
            .attributes = Attribute::Bold,
        },
    },
};
\end{lstlisting}

No allocation: \code{string\_view} into the literal, small POD
style. The whole construction is on the stack.

For a box:

\begin{lstlisting}
Element panel{
    .v = BoxElement{
        .children = {
            Element{.v = TextElement{.text = std::string_view{"Top"}}},
            Element{.v = TextElement{.text = std::string_view{"Mid"}}},
            Element{.v = TextElement{.text = std::string_view{"Bot"}}},
        },
        .layout = FlexStyle{.direction = Direction::Column},
        .border = Border{.style = BorderStyle::Rounded},
        .padding = Padding::all(1),
    },
};
\end{lstlisting}

The \code{children} list contains nested \code{Element} values.
One allocation for the vector buffer; the children themselves are
inline within their slots.

\subsection*{The box() builder helper}

For tests and one-off scripts, \code{box(...)} bundles common
defaults:

\begin{lstlisting}
auto e = box(text("hi"))
             .border(BorderStyle::Single)
             .padding(1)
             .build();
\end{lstlisting}

The builder is a small fluent API that mutates an internal
\code{BoxElement} and returns the wrapped \code{Element} on
\code{build()}. Use it when designated initialisers feel verbose.

\section{The colour model, up close}

\code{Color} is a sum of three representations:

\begin{lstlisting}
struct Color {
    enum class Kind { Default, Ansi16, Ansi256, Rgb };
    Kind kind = Kind::Default;
    uint8_t r = 0, g = 0, b = 0;  // payload for Rgb
    uint8_t index = 0;            // payload for Ansi16/Ansi256
    bool is_default() const { return kind == Kind::Default; }
};
\end{lstlisting}

The four kinds:

\begin{itemize}[leftmargin=*]
  \item \textbf{Default.} ``No colour set; inherit from terminal.''
    The ANSI emitter writes nothing for this case.
  \item \textbf{Ansi16.} Indexes 0--15 from the classic palette
    (black, red, green, \ldots, plus bright variants). Maps to
    SGR codes 30--37 / 40--47 and 90--97 / 100--107.
  \item \textbf{Ansi256.} Indexes 0--255 from the 256-colour
    palette (16 ANSI + 6$\times$6$\times$6 colour cube + 24-step
    grey ramp). Maps to SGR 38;5;$n$ / 48;5;$n$.
  \item \textbf{Rgb.} 24-bit truecolour. Maps to
    SGR 38;2;$r$;$g$;$b$ / 48;2;$r$;$g$;$b$.
\end{itemize}

\subsection*{Why Default exists}

``No colour'' is not the same as ``black''. If the user has a
dark-on-light terminal theme, drawing ``black'' over their
background produces invisible text or a dark blob. Drawing ``no
colour'' lets the terminal supply the foreground, which adapts
to the user's theme.

The DSL exposes \code{Color::default\_()} (or, in the runtime
builders, leaves the \code{Color} field default-constructed) to
opt out of overriding the user's choice.

\subsection*{8/256/truecolour: a small history}

\begin{itemize}[leftmargin=*]
  \item The original ANSI standard (X3.64, 1979) defined eight
    colours: SGR 30--37.
  \item The bright variants (90--97) were added by \emph{aixterm}
    around 1995 and adopted widely.
  \item The 256-colour palette was popularised by xterm in 1999
    via the SGR 38;5;$n$ syntax; the cube and grey ramp were
    xterm's invention, not a standard.
  \item Truecolour (24-bit RGB) was first implemented by
    Konsole and iTerm2 in the early 2010s; ITU-T T.416 had
    described the syntax \code{38:2::r:g:b} since 1993, but
    everyone uses \code{38;2;r;g;b} now.
\end{itemize}

\begin{funfact}
If you read T.416 carefully, the separators between the RGB
values should be colons (\code{:}), not semicolons (\code{;}).
Because every implementation used semicolons, the de-facto
standard split from the de-jure one. \maya{} emits semicolons,
like everyone else.
\end{funfact}

\subsection*{Choosing a colour kind}

When you write \code{Fg<R, G, B>} in the DSL, you get
\code{Color::Kind::Rgb}. When you write \code{Fg<"red">} you get
\code{Color::Kind::Ansi16}. The runtime emitter degrades
automatically: if the user's terminal does not advertise
truecolour (\code{TERM} or \code{COLORTERM} environment), the
emitter quantises RGB to the nearest 256-colour index.

The quantisation uses a simple Euclidean-distance match against
the 256-colour palette. The result is usually visually close;
for brand-perfect colour, set \code{COLORTERM=truecolor} in your
shell config.

\section{The border gallery}

\code{BorderStyle} is an enum. Each value picks a set of six
byte sequences (corners, horizontal, vertical) drawn around the
box. The Unicode block-drawing range
U+2500--U+257F supplies all the glyphs.

\subsection*{Single (U+250C--U+2518, U+2500, U+2502)}

The canonical thin-line border:

\begin{terminalbox}
+--Border-------+\\
|\, body line  \, |\\
+--------------+
\end{terminalbox}

Corners use U+250C (top-left), U+2510 (top-right), U+2514 (bottom-
left), U+2518 (bottom-right). Horizontals use U+2500; verticals
use U+2502.

\subsection*{Rounded (U+256D--U+2570)}

The friendly variant with curved corners:

\begin{terminalbox}
.--Rounded----.\\
|\, body line\, |\\
`-------------'
\end{terminalbox}

Corners are U+256D / U+256E / U+2570 / U+256F. Lines reuse
U+2500 and U+2502.

\subsection*{Double (U+2554, U+2557, U+255A, U+255D, U+2550, U+2551)}

Line-drawing for headers and important regions:

\begin{terminalbox}
=====================\\
||\, body line\, ||\\
=====================
\end{terminalbox}

(Real characters: U+2554 / U+2557 / U+255A / U+255D corners,
U+2550 horizontal, U+2551 vertical.)

\subsection*{Heavy (U+250F, U+2513, U+2517, U+251B, U+2501, U+2503)}

A bold thick-line border for emphasis.

\subsection*{ASCII}

A portability fallback using only ASCII: hyphens, pipes, and
plus signs. Use this on terminals that cannot render Unicode
block-drawing characters (rare, but they exist).

\subsection*{Custom}

A \code{Border} struct can also carry six \code{char32\_t}
values directly, bypassing the enum. Useful for half-block fills
and custom themes.

\subsection*{Partial-side borders}

\code{Border} carries a \code{sides} bitmask (top, right, bottom,
left). Set only some bits to get a border on a subset of sides:

\begin{lstlisting}
Border b{
    .style = BorderStyle::Single,
    .sides = BorderSide::Bottom | BorderSide::Top,
};
\end{lstlisting}

A box with only top and bottom borders looks like a horizontal
rule sandwich --- useful for table rows and section dividers.

\section{Style: merging rules}

A \code{Style} struct carries:

\begin{itemize}[leftmargin=*]
  \item \code{fg}, \code{bg}: \code{Color} values (or default).
  \item \code{attributes}: a bitset of
    Bold/Dim/Italic/Underline/Blink/Reverse/Strike.
  \item \code{hyperlink}: optional URL for OSC 8.
\end{itemize}

When styles compose (a child inherits from a parent, or two
pipes attach styles), \maya{} applies a merge rule:

\begin{enumerate}[leftmargin=*]
  \item Attributes union: if either has Bold, the result has
    Bold.
  \item Colours override: a non-default colour replaces a
    default. Two non-default colours: child wins.
  \item Hyperlink: child wins if set.
\end{enumerate}

\begin{lstlisting}
Style merge(Style outer, Style inner) {
    Style r = outer;
    r.attributes |= inner.attributes;
    if (!inner.fg.is_default()) r.fg = inner.fg;
    if (!inner.bg.is_default()) r.bg = inner.bg;
    if (inner.hyperlink)        r.hyperlink = inner.hyperlink;
    return r;
}
\end{lstlisting}

The merge happens during the layout pass, when the renderer walks
the tree and assembles per-cell style. You do not call it
yourself; \maya{} does.

\begin{idea}
The merge rule is what makes
\code{box(text("x")|Italic)|Bold} produce text that is both bold
and italic, without you writing \code{Bold | Italic} explicitly.
Styles inherit through the tree the way CSS does, just with a
cleaner rule set.
\end{idea}

\section{Tree-walking patterns}

Once you have an \code{Element}, every interesting operation is a
recursive visit. Five patterns appear over and over.

\subsection*{Counting nodes of a kind}

\begin{lstlisting}
int count_text(const Element& e) {
    return std::visit(overload{
        [](const TextElement&)      { return 1; },
        [](const BoxElement& b) {
            int n = 0;
            for (auto& c : b.children) n += count_text(c);
            return n;
        },
        [](const ElementList& l) {
            int n = 0;
            for (auto& c : l.items) n += count_text(c);
            return n;
        }
    }, e.v);
}
\end{lstlisting}

Visit at the root, recurse into containers. The visit dispatches
on kind; the recursion threads through the children.

\subsection*{Finding the first text matching a predicate}

\begin{lstlisting}
std::optional<std::string_view>
find_text(const Element& e,
          std::function<bool(std::string_view)> pred) {
    return std::visit(overload{
        [&](const TextElement& t) -> std::optional<std::string_view> {
            auto sv = std::visit(
                [](auto& x) { return std::string_view{x}; }, t.text);
            return pred(sv) ? std::optional{sv} : std::nullopt;
        },
        [&](const BoxElement& b) -> std::optional<std::string_view> {
            for (auto& c : b.children)
                if (auto r = find_text(c, pred)) return r;
            return std::nullopt;
        },
        [&](const ElementList& l) -> std::optional<std::string_view> {
            for (auto& c : l.items)
                if (auto r = find_text(c, pred)) return r;
            return std::nullopt;
        }
    }, e.v);
}
\end{lstlisting}

Classic depth-first search. The optional return is the standard
way to express ``found / not found'' without exceptions.

\subsection*{Transforming text in place}

\begin{lstlisting}
void uppercase_text(Element& e) {
    std::visit(overload{
        [](TextElement& t) {
            std::string s;
            std::visit([&](auto& x) { s = std::string{x}; }, t.text);
            for (auto& c : s) c = std::toupper(c);
            t.text = std::move(s);
        },
        [](BoxElement& b) {
            for (auto& c : b.children) uppercase_text(c);
        },
        [](ElementList& l) {
            for (auto& c : l.items) uppercase_text(c);
        }
    }, e.v);
}
\end{lstlisting}

Mutating visitor: the lambdas take non-const references and
modify in place. The variant remains the same kind; only its
payload changes.

\subsection*{Hit-testing (mouse coordinate to element)}

Given a screen coordinate \code{(x, y)}, find the deepest
element under that point. Layout has annotated each box with its
bounding rectangle:

\begin{lstlisting}
const Element* hit_test(const Element& e, int x, int y) {
    return std::visit(overload{
        [&](const TextElement& t) -> const Element* {
            // text knows its rect; check membership
            return t.rect.contains(x, y) ? &e : nullptr;
        },
        [&](const BoxElement& b) -> const Element* {
            if (!b.rect.contains(x, y)) return nullptr;
            // Recurse to find a deeper hit, else return self.
            for (auto& c : b.children)
                if (auto h = hit_test(c, x, y)) return h;
            return &e;
        },
        [&](const ElementList& l) -> const Element* {
            for (auto& c : l.items)
                if (auto h = hit_test(c, x, y)) return h;
            return nullptr;
        }
    }, e.v);
}
\end{lstlisting}

Returns the deepest hit, or \code{nullptr} if the coordinate is
outside the entire tree.

\subsection*{Pretty-printing for debugging}

\begin{lstlisting}
void dump(const Element& e, int depth = 0) {
    auto pad = std::string(depth * 2, ' ');
    std::visit(overload{
        [&](const TextElement& t) {
            std::visit([&](auto& s) {
                std::cout << pad << "text: " << s << "\n";
            }, t.text);
        },
        [&](const BoxElement& b) {
            std::cout << pad << "box (" << b.children.size() << " children)\n";
            for (auto& c : b.children) dump(c, depth + 1);
        },
        [&](const ElementList& l) {
            std::cout << pad << "list (" << l.items.size() << ")\n";
            for (auto& c : l.items) dump(c, depth + 1);
        }
    }, e.v);
}
\end{lstlisting}

Useful while you are learning the tree shape. Drop a
\code{dump(my\_element)} in a test and you can see what the DSL
produced.

\begin{recap}
Every operation on an \code{Element} is a recursive
\code{std::visit}. The pattern is uniform: handle each
alternative, recurse into containers. There is no virtual
dispatch; the compiler turns the visit into a switch.
\end{recap}

\section{Hyperlinks: OSC 8}

Xterm and most modern terminals support clickable hyperlinks via
the OSC 8 escape sequence:

\begin{lstlisting}
\e]8;;https://example.com\e\\Visit\e]8;;\e\\
\end{lstlisting}

The text between the start (\code{ESC ]8;;URL ESC \textbackslash})
and end (\code{ESC ]8;; ESC \textbackslash}) sequences is
rendered as a clickable region.

\maya{}'s \code{TextElement::hyperlink} field, if set, causes the
emitter to wrap the cell in OSC 8 brackets. The text itself is
unchanged; only the metadata differs.

\begin{lstlisting}
Element link{
    .v = TextElement{
        .text = std::string_view{"the docs"},
        .style = Style{.attributes = Attribute::Underline},
        .hyperlink = std::string{"https://example.com"},
    },
};
\end{lstlisting}

Rendering this in a terminal that supports OSC 8 produces a
clickable ``the docs''. On a terminal that does not, the
emitter falls back to plain text (the hyperlink sequence is
ignored).

\begin{aside}
OSC 8 was first implemented by VTE (the terminal widget in
GNOME Terminal) in 2017 and is now supported by iTerm2, Kitty,
WezTerm, Alacritty, and modern xterm. Tmux passes it through
as of version 3.3. If your terminal is older than
2019-ish, expect plain text.
\end{aside}

\section{Pitfalls}

A few traps to watch for when working with elements directly.

\begin{warning}
\textbf{Mixing string\_view and string carelessly.} If a
\code{TextElement::text} holds a \code{string\_view} that points
into a temporary, the view will dangle. Use \code{string\_view}
only for literals and other long-lived storage; use
\code{std::string} for anything computed.
\end{warning}

\begin{warning}
\textbf{Copying large element trees.} A deep tree's copy is
proportional to its depth and width. If you find yourself
cloning elements every frame, you are paying for moves that
should be moves and copies that should be views. Profile
before you optimise --- the cost is rarely the bottleneck ---
but watch the allocator.
\end{warning}

\begin{warning}
\textbf{Setting Visibility::Collapsed and expecting layout to
stay stable.} Collapsed removes the element from layout
entirely; its siblings will reflow. If you want a stable layout,
use \code{Hidden} instead, which reserves space.
\end{warning}

\begin{warning}
\textbf{Custom Border characters with combining marks.} A single
Unicode codepoint can take zero, one, or two cells. \maya{}'s
border drawing assumes one cell per character. If you supply a
combining mark (e.g. U+0301), the border will misalign. Stick to
the block-drawing range.
\end{warning}

\begin{warning}
\textbf{Hyperlinks in text that wraps.} If a hyperlink text
wraps across two visual lines, some terminals split the link
into two regions; some keep them unified. The behaviour is
terminal-dependent; do not rely on either.
\end{warning}

\section{The bigger picture}

The element tree is what flows between \code{view} and the
renderer. Every frame, \code{view} returns a fresh \code{Element};
the renderer walks it, lays it out, diffs it, and emits bytes.
The tree is the lingua franca between your code and \maya{}'s
machinery.

Understanding the tree pays off in three places:

\begin{itemize}[leftmargin=*]
  \item Debugging a misrendered frame: \code{dump} the tree and
    spot the wrong child.
  \item Writing widget libraries: a widget is a function from
    state to \code{Element}, so you must build trees by hand.
  \item Reading \maya{}'s source: the layout, the diff, the
    ANSI emitter --- all walk the same variant.
\end{itemize}

In the next chapter we follow the tree into the layout engine
and see how the abstract description becomes a grid of cells.

\begin{recap}
An \code{Element} is a closed sum: \code{BoxElement},
\code{TextElement}, \code{ElementList}. Three small structs, one
variant. Every \maya{} algorithm visits it. There is no
inheritance; the set of node kinds is fixed and exhaustively
checkable. Constructing trees by hand is one designated
initialiser at a time.
\end{recap}
